% !BIB program = biber 
\documentclass{article}

%----------------------------------------------------------------------------------------
%	PACKAGES AND OTHER DOCUMENT CONFIGURATIONS
%----------------------------------------------------------------------------------------

% Generate dummy text (lorem ipsum) in your document
\usepackage{lipsum} % Comment me

% English language hyphenation
\usepackage[english]{babel}

% Math packages
\usepackage{amsmath, amsfonts, amsthm}

% Indent first paragraph
\usepackage{indentfirst}

% Figures
\usepackage{float}
\usepackage{graphicx}
\graphicspath{{./images/}}

% Captions
\usepackage{caption}
\usepackage{subcaption}

% Context sensitive quotation facilities
\usepackage{csquotes}

% Enumerate
\usepackage{enumitem}

% Chemical formulas package
\usepackage{chemformula}

% Filesystem
\usepackage{forest}

% Minted package for syntax highlighting
\usepackage{minted}

% Table 
\usepackage[toc]{appendix}
\addto\captionsenglish{% Replace "english" with the language you use
  \renewcommand{\contentsname}{Table of Contents}
}

% Shorthands
\usepackage{xspace}
\makeatletter
\DeclareRobustCommand\onedot{\futurelet\@let@token\@onedot}
\def\@onedot{\ifx\@let@token.\else.\null\fi\xspace}

\def\eg{e.g\onedot} \def\Eg{E.g\onedot}
\def\ie{i.e\onedot} \def\Ie{I.e\onedot}
\def\cf{cf\onedot} \def\Cf{Cf\onedot}
\def\etc{etc\onedot} \def\vs{vs\onedot}
\def\wrt{w.r.t\onedot} \def\dof{d.o.f\onedot}
\def\etal{et al.\onedot}
\makeatother

% Hyperlinks styling
\usepackage{hyperref}
% Hyperlink setup
\hypersetup{
  colorlinks   = true, % Colours links instead of ugly boxes
  urlcolor     = blue, % Colour for external hyperlinks
  linkcolor    = blue, % Colour of internal links
  citecolor    = blue, % Colour of citations
  linktocpage  = true, % Link only on page numbers
}
% PDF metadata
\hypersetup{
  pdftitle     = {Guia do usuário: Co-localização de parasitas de Chagas em
      células musculares cardíacas},
  pdfauthor    = {Patrick. H. F. Alvares, João V. S. Guerra, José G. C.
      Pereira},
  pdfproducer  = {Latex with hyperref},
  pdfcreator   = {pdflatex}
}

% Bibliography
\usepackage[
  backend=biber,
  style=ieee,
  natbib=true,
]{biblatex}
\addbibresource{\jobname.bib}

% Acronyms
\usepackage[
  printonlyused,
  nohyperlinks,
  withpage,
]{acronym}

%----------------------------------------------------------------------------------------
%	DOCUMENT MARGINS
%----------------------------------------------------------------------------------------

\usepackage{geometry} % Required for adjusting page dimensions and margins
\geometry{
  paper=a4paper,
  top=30mm, % Top margin
  bottom=20mm, % Bottom margin
  left=30mm, % Left margin
  right=20mm, % Right margin
  % showframe, % Uncomment to show how the type block is set on the page
}

%----------------------------------------------------------------------------------------
%	FONTS
%----------------------------------------------------------------------------------------

% Required for inputting international characters
\usepackage[T1]{fontenc}
% Use 8-bit encoding
\usepackage[latin1,utf8]{inputenc}

% Fonts
\usepackage{lmodern}

% Title Page
\long\def\reporttitle{%
  \begin{titlepage}
    % Logo
    \begin{center}
      \begin{minipage}{0.9\textwidth}%

        \includegraphics[width=\textwidth]{lnbio-cnpem-logo.png}%
      \end{minipage}
    \end{center}
    \vskip 0.5em

    % Horizontal line
    \hline
    \vskip 1em

      % Title
      {\centering \LARGE Guia do usuário: Co-localização de parasitas de
        Chagas em células musculares cardíacas \par}
    \vskip 1em

      % Author
      {\centering \large Patrick. H. F. Alvares, João V. S. Guerra, José G.
        C. Pereira \par} % CHANGE
    \vskip 1em

      % Date
      {\centering \large \today\par}
    \vskip 0.5em

    % Horizontal line
    \hline

    % Do not break page
    \let\endtitlepage\relax
  \end{titlepage}
}

\begin{document}

% Make title
\reporttitle
\break

% Table of Contents
\hline
\tableofcontents
\vskip 0.5em
\hline
\break 

\section{Introdução}

Este guia do usuário fornece instruções detalhadas sobre como executar a análise de co-localização de parasitas de Chagas em células musculares cardíacas. O objetivo deste guia é auxiliar os usuários a configurar e executar o pipeline de análise, bem como a interpretar os resultados obtidos.

\section{Requisitos}

Para executar a análise de co-localização de parasitas de Chagas em células musculares cardíacas, você precisará dos seguintes requisitos:

\begin{itemize}
  \item \textbf{Acesso ao GitHub do CNPEM}: Você deve ter acesso ao repositório do projeto (\url{https://github.com/cnpem/ParasiteCoLocalization}) no GitHub do CNPEM (\url{https://github.com/cnpem/}) para baixar os arquivos necessários. Se não tiver acesso, entre em contato com a Equipe de Dados Biológicos (\url{edb@lnbio.cnpem.br}).
  \item \textbf{Acesso ao HPC Marvin}: Você deve ter acesso ao HPC Marvin para executar o pipeline de análise. Se você não tiver acesso, acesse \url{https://marvindocs.cnpem.br/primeiros-passos/index.html}.
\end{itemize}

\section{Acessando o HPCC Marvin}

Para acessar o HPCC Marvin, comece abrindo o terminal. Se estiver usando Windows, abra o PowerShell; se estiver usando Linux ou MacOS, abra o terminal Terminal. Para acessar o HPCC Marvin, use o comando:

\begin{minted}[frame=single, bgcolor=lightgray, fontsize=\small]{bash}
ssh <seu.login.cnpem>@marvin.cnpem.br
\end{minted}

Após inserir o comando, será solicitado que você insira sua senha, digite sua senha institucional.

\section{Clonando o Repositório ParasiteCoLocalization}

Para clonar o repositório \href{https://github.com/cnpem/ParasiteCoLocalization}{ParasiteCoLocalization}, execute o seguinte comando no terminal:

\begin{minted}[frame=single, bgcolor=lightgray, fontsize=\small]{bash}
git clone https://github.com/cnpem/ParasiteCoLocalization.git
\end{minted}

O comando acima cria uma pasta chamada \texttt{ParasiteCoLocalization} no diretório atual. Esta pasta contém todos os arquivos necessários para a execução da análise de co-localização de parasitas de Chagas em células musculares cardíacas, incluindo o script \texttt{run.sh}, que automatiza a execução da análise via SLURM.

\subsection{Organização do repositório \texttt{ParasiteCoLocalization}}

A estrutura de diretórios do projeto é apresentada na árvore de arquivos a seguir:

\begin{forest}
  for tree={
  font=\ttfamily,
  grow'=0,
  child anchor=west,
  parent anchor=south,
  anchor=west,
  calign=first,
  edge path={%
      \noexpand\path [draw, thick, \forestoption{edge}]
      (!u.south west) |- (.child anchor)\forestoption{edge label};
    },
  before typesetting nodes={
      if n=1
        {insert before={[, phantom]}}
        {}
    },
  fit=band,
  before computing xy={l=15pt},
  }
  [./ParasiteCoLocalization/
  [benchmarking]
  [data/]
  [docs/]
  ]
\end{forest}

\section{Enviando Imagens para o HPC Marvin}

Para enviar as imagens para o HPC Marvin, você pode usar o comando \texttt{scp} no terminal. O comando \texttt{scp} permite copiar arquivos de e para servidores remotos de forma segura.

## Preencher

Para enviar um arquivo para o HPC Marvin, use o seguinte comando:

## Preencher

\section{Executando o Pipeline no HPC Marvin}

Após enviar as imagens para o HPC Marvin, você pode executar o pipeline de análise de co-localização de parasitas de Chagas em células musculares cardíacas. O script \texttt{run.sh} executa o pipeline de análise via SLURM, que é o sistema de gerenciamento de tarefas usado no HPC Marvin. Para isso, execute o script \texttt{run.sh} no diretório \texttt{ParasiteCoLocalization}.

\begin{minted}[frame=single, bgcolor=lightgray, fontsize=\small]{bash}
cd ParasiteCoLocalization
sbatch run.sh -m marvin
\end{minted}

O comando acima retorna o ID do job SLURM, que você pode usar para monitorar o progresso da análise.

\subsection{Monitoramento do Job SLURM}

Para monitorar o progresso do job SLURM, você pode usar o comando \texttt{squeue} no terminal. O comando \texttt{squeue} exibe informações sobre os jobs SLURM em execução no HPC Marvin.

\begin{minted}[frame=single, bgcolor=lightgray, fontsize=\small]{bash}
squeue -u <seu.login.cnpem>
\end{minted}

Se o job estiver em execução, o estado será exibido como R (em execução). Se o job estiver aguardando na fila, o estado será PD (pendente).

Caso deseje cancelar o job SLURM, você pode usar o comando \texttt{scancel} no terminal. O comando \texttt{scancel} cancela um job SLURM em execução no HPC Marvin.

\begin{minted}[frame=single, bgcolor=lightgray, fontsize=\small]{bash}
scancel <job_id>
\end{minted}

Por fim, você pode verificar o arquivo de log gerado pelo SLURM (\texttt{slurm\_<job\_id>.out}) para obter informações sobre o seu job.

\section{Glossário}

Here is an \ac{Ex} of acronym usage.

\begin{acronym} \itemsep=-8pt
  \acro{Ex}[Ex]{Example}
\end{acronym}

\printbibliography

% \appendix

% \break \section{Supplementary Information}

% \renewcommand{\thefigure}{S\arabic{figure}}
% \setcounter{figure}{0}
% \renewcommand{\thetable}{S\arabic{table}}
% \setcounter{table}{0}

% \begin{table}[H]
%     \begin{center}
%         \begin{tabular}{|c|c|}
%             \hline
%             a & b \\\hline
%             c & d \\\hline
%         \end{tabular}
%     \end{center}
%     \caption[Shorter table caption]{Table caption caption caption caption
%         caption caption caption caption caption caption caption caption caption
%         caption caption caption caption caption caption caption caption caption
%         caption caption.}
%     \label{t:label0}
% \end{table}

% \begin{figure}[H]
% \centerline{\includegraphics[scale=0.2]{lnbio-cnpem-logo.png}}
% \caption[Shorter figure caption]{Figure Caption caption caption caption
%   caption caption caption caption caption caption caption caption caption
%   caption caption caption caption caption caption caption caption caption
%   caption caption.}
% \label{f:label1}
% \end{figure}

\end{document}
